% 本文件由 LaTeX 排版 Agent 根据有效 translation.json 自动生成。
% author=Tertullian; work_id=0406; title=论一夫一妻制

\chapter{论一夫一妻制}

% structure: p0001 source=h1 target=omitted reason=page-title-already-rendered-by-parent

% structure: p0002 source=h3 target=section reason=relative-heading-rank; base=h3

\section{第一章 异端者、属魂者与属神者所持的不同婚姻观}

异端者废除婚姻；属魂者则积累婚姻。前者甚至连一次也不结婚；后者不止一次。造物主的法律啊，你该怎么办？在异族阉人与自家的仆人之间，你既抱怨自家的仆人过分顺服，又抱怨外人的轻蔑。滥用你的人与不使用你的人，同样伤害了你。事实上，这样的节制并不可赞，因为它是异端的；这样的放纵也不可辩护，因为它是属魂的。前者亵渎，后者放荡；前者摧毁婚姻之天主，后者使祂蒙羞。然而，在我们中间——因承认属神恩赐而堪当被称为属神者——节制之虔诚如同放纵之谦逊；因为二者都与造物主和谐。节制尊重婚姻的法律，放纵则调节它；前者并非被迫，后者受到规范；前者承认自由意志的能力，后者承认一种限制。我们只承认一次婚姻，如同我们只承认一位天主。婚姻的法律因与廉耻相联而获得更大的尊荣。但属魂者既未领受圣神，圣神的事就不为他们所喜悦。因此，只要圣神的事不令他们喜悦，血肉的事就会令他们喜悦，因为它们是圣神的对立面。宗徒说：\BibleQuote{血肉相反圣神，圣神相反血肉。}但血肉会\Quote{贪恋}什么呢？除非是更属于血肉的事。为此，在起初，血肉就与圣神疏远了。天主说：\BibleQuote{我的神不能永远留在这人身上，因为他们是血肉。}（\BibleRef{创 6:3}）\par

% structure: p0004 source=h3 target=section reason=relative-heading-rank; base=h3

\section{第二章 属神者从标新立异的指控中得到辩护}

于是，他们指责一夫一妻的纪律是异端；也没有别的理由使他们不得不否认护慰者，而是因为他们认为祂创立了一种新的纪律，一种他们觉得极其苛刻的纪律：因此，这是我们首先必须从一般处理（主题）上切入的问题：是否有理由坚持护慰者所教导的任何事，可以被指责为（在反对大公传统的意义上）标新立异，或（在反对主的\BibleQuote{轻省的轭}（\BibleRef{玛 11:30}）的意义上）沉重？\par

关于这两点，主自己已经宣判。因为当祂说：\BibleQuote{我还有许多事要告诉你们，但你们现在不能承受。当圣神来了，祂要引导你们进入一切真理}（\BibleRef{若 16:12-13}），祂充分地向我们表明，祂将带来那些既被视为\Emph{新颖}（因为从未被宣布过）又最终被视为\Emph{沉重}（仿佛那就是未被宣布的原因）的教导。你说：\Quote{由此推论，任何你喜欢的、新颖而沉重的事，都可以归给护慰者，即使它来自敌对神。}当然不。因为敌对神会因祂教导的差异性而显明出来，首先通过篡改信仰准则，然后（继续）篡改纪律的秩序；因为那占据首要级别者（即先于纪律的信仰）的败坏，是首先发生的。一个人必然首先对天主持异端看法，然后才对其制度。但护慰者，有许多事要完全教导，这些事主推迟到祂来时才说（按照预定的安排），祂将首先为基督作强有力的见证，正如我们所相信的（祂是）那样，连同整个造物主天主的秩序，并将荣耀祂，并\BibleQuote{使祂\Person{耶稣}的一切话}（\BibleRef{若 14:26}）为人\Quote{记起}。当祂根据基本原则被认出来时，祂将揭示那些属于纪律的\BibleQuote{许多事}（\BibleRef{若 16:12}）；而祂教导的完整无缺为其（启示）赢得信任，尽管它们\Quote{新颖}（因为它们现在才被揭示），尽管它们\Quote{沉重}（因为甚至\Emph{现在}仍被认为难以承受）：但这些启示所涉及的，正是那位曾说祂还有\BibleQuote{许多事}（\BibleRef{若 16:12}）要由护慰者完全教导的基督，这些事对当时代的人同样沉重，正如对那些人（他们当时还\Quote{不能承受}）一样沉重。\par

% structure: p0007 source=h3 target=section reason=relative-heading-rank; base=h3

\section{第三章 结合主与宗徒之言进一步考量标新立异的问题}

但是，（至于）独身是否\Quote{难以负担}，让那仍然无耻的\Quote{肉体的软弱}来考虑吧：我们暂且先达成一致，看它是否\Quote{新奇}。我们甚至作出这一更宽泛的断言：即使护慰者在这我们的时代明确规定了完全的、绝对的童贞或节欲，以至于不允许肉体的情欲甚至在单身婚姻中发泄出来，即便如此，祂似乎也没有引入任何\Quote{新奇}的东西；因为主亲自为\Quote{阉人}开启了\Quote{诸天的国度}，因为祂自己也是童贞者；宗徒也仰望祂——他自己也因此而禁欲——更推崇节欲。（\Quote{是的}，）你说，\Quote{但保留婚姻的法律。}显然保留，我们将看在何种限制下；然而，祂已经在某种程度上摧毁了它，因为他更推崇节欲。他说：\Quote{男人不接触女人是好的。}因此，接触她便是恶的；因为除了恶，没有什么是与善对立的。因此（他说）：\Quote{剩下的就是，那些有妻子的，要像没有一样}，这样对没有妻子的人来说，避免拥有妻子就更有约束力。他也给出了这样劝告的理由：未婚的人思虑天主的事，但已婚的人思虑如何在婚姻中取悦对方。我可以争辩，那被\Emph{允许}的并非绝对善。因为绝对善不是被\Emph{允许}的，而是不需要请求使其合法。许可有时甚至出于\Emph{必要}。最后，在这种情况下，许可婚姻的一方并没有意愿。因为他的\Emph{意愿}指向别处。他说：\Quote{我\Emph{愿意}你们都像我一样。}当他表明（如此保持）是\Quote{更好的}，请问，他表明自己\Quote{愿意}什么呢，岂非他之前所预设的\Quote{更好的}？因此，如果他\Emph{允许}某些与他所\Quote{愿意}的不同——并非自愿允许，而是出于必要——他表明他非自愿所宽许的并非绝对善。最后，当他说：\Quote{结婚比燃烧更好}，我们该如何理解这种善呢，它比一种惩罚更好？只有当与非常坏的事物比较时，它才能显得\Quote{更好}。\Quote{善}是自身保持此名者；无需比较——我不是说与恶比较，甚至——与别的善比较：这样，即使与另一种善比较并被遮蔽，它仍然保有善的名称。另一方面，与恶的比较是迫使它被称为善的手段；它与其说是\Quote{善}，不如说是一种较低级的恶，当被更高的恶遮蔽时，便被逼称为善。简而言之，去掉条件，不说\Quote{结婚比燃烧更好}；我怀疑你是否敢于说\Quote{（结婚）更好}，而不加上比什么更好。这样做之后，它就不再是\Quote{更好的}；同时，既不是\Quote{更好的}，也不是\Quote{善的}，因为条件被去掉，而这条件使它在与另一事物比较时成为\Quote{更好的}，从而迫使它被视为\Quote{善}。失去一只眼睛比失去两只更好。然而，如果你取消两者的比较，那么失去一只眼睛便不是更好的，因为它甚至不是善的。\par

那么，如果他出于我们提到的必要性，基于他自己的（即人的）感觉，通融地准许一切结婚的宽容，因为\Quote{结婚比燃烧更好}？事实上，当他转向第二种情况，说：\Quote{至于已婚者，我正式宣布——不是我，而是主}——他表明他上面所说的并非主的权威，而是人的判断。然而，当他引导他们回心转意趋向节欲时，（\Quote{但我愿意你们都如此}），他说：\Quote{而且我想，我也有天主的圣神}；这样，如果他出于必要而准许了某种宽容，他可以藉圣神的权威召回。但是，若望在劝告我们\Quote{应当像主那样行事}时，当然也提醒我们按肉体的圣洁（如同按祂在其他方面的榜样）行事。因此，他更明确地说：\Quote{凡对祂怀有这希望的，就使自己纯洁，如同祂是纯洁的一样。}因为别处又（读到）：\Quote{你们要圣洁，如同祂是圣洁的一样}——即在肉体上。因为论到圣神，他不会那样说，因为圣神无需任何外在影响就被认为是\Quote{圣洁的}，也不等待被劝告要圣洁，这是祂的本性。但肉体是\Emph{被教导}圣洁的；而且在基督身上，它也是圣洁的。\par

因此，如果所有这些（考虑）废除了婚姻的许可，无论我们看许可的条件，还是强加的节欲的优先性，为什么在宗徒之后，同一位圣神，为了引导门徒身份进入\Quote{一切真理}而随着时代逐步（按照传道者所说，\Quote{凡事都有定时}），到此刻还不给肉体加上最终的笼头，不再含蓄地叫我们远离婚姻，而是公开地？因为如今（比以往更）\Quote{时候已经缩短了}——从那时起约一百六十年过去了？你不会不自觉地（这样）思考吗：\Quote{这纪律是古老的，早已在主的肉身和意志中预先显示，随后在祂宗徒的劝告和榜样中相继显现？我们自古以来就被预定这圣洁。护慰者没有引入任何新奇。祂所预言的，祂（现在）最终指定；祂所推迟的，祂（现在）要求。}紧接着，思索这些，你会轻易说服自己，护慰者更有资格传讲一夫一妻制，祂本也可以传讲废除婚姻；而且更可信的是祂缓和了祂本可废除的事，如果你明白基督的\Quote{意志}是什么。在此你也应当认出护慰者作为安慰者的品格，因为祂以绝对的节欲（的严格）来原谅你的软弱。\par

% structure: p0011 source=h3 target=section reason=relative-heading-rank; base=h3

\section{第四章 略过护慰者，\Person{特土良}转向古代圣经及其关于当前主题的见证}

现在，暂且不提护慰者作为我们自己的某种权威，我们来展开原始圣经的通用工具。这件事是我们可以证明的：一夫一妻制的规则既不是新奇的，也不是奇怪的，反而既是古老的，又是基督徒所特有的；以便你意识到护慰者是其\Emph{恢}复者，而不是\Emph{创}立者。至于年代久远，还有什么比人类原始的源头本身更古老的正式模式可以提出呢？天主为男人造了一个女人，从他身上取了一根肋骨，并且（当然）是从众多中取出一根。但是，在造物工作之前的开场白中，祂说：\Quote{那人独居不好；让我们为他造一个相称的助手吧。}因为如果祂预定他（男人）要有多个妻子，祂就会说\Quote{助手们}。祂还加了一条关于未来的法律；也就是说，（那些话）\Quote{二人要成为一体}——不是三人，也不是更多；否则，若更多，他们就不再是\Quote{二人}——是预言性地宣告的。这法律（坚定）屹立。简言之，婚姻的合一在我们种族祖先那里一直持续到底；不是因为当时没有其他女人，而是因为没有其他女人的原因\Emph{为何}，是为了不让初熟的果实被双重婚姻玷污。否则，如果天主（如果）愿意，祂本来\Emph{可以}有（有其他女人）；无论如何，祂可以从自己女儿的丰裕中取——既然厄娃是从他自己的骨肉中取出的——如果敬虔允许这样做的话。但是，哪里有第一桩罪行（即）杀人，以谋杀兄弟开始，就没有什么罪行比双重婚姻更配得第二的位置。因为一个人先后有两个妻子，或者同时有两个（妻子）使得（数目为）二，这两者并无区别。结合的和分离的（个体）数目相同。尽管如此，天主的制度，在通过\Person{拉默客}受到一次暴力侵犯之后，在那种族中一直保持稳固到最后。没有出现第二个\Person{拉默客}，以拥有两个妻子的丈夫的方式。圣经不提的，就是它否认的。其他不义引发了洪水：（不义）一次受罚，无论其性质如何；然而，不是\Quote{七十七次}，那是双重婚姻所应得的（惩罚）。\par

但是再次：第二次人种的改革是从一夫一妻制为其母追溯的。再次，\Quote{二人成为一体}承担了\Quote{生养众多}的（职责）——\Person{诺厄}（即）和他的妻子，以及他们的儿子，在单一婚姻中。甚至在动物身上也承认一夫一妻制，以免连走兽都从奸淫而生。\Quote{从所有走兽中，}（天主）说，\Quote{从一切血肉中，你要带你进入方舟，使它们与你一同存活，公母成对：它们要从所有飞禽按（它们的）种类，以及地上一切爬虫按它们的种类中取来；所有之中要进入你那里，公母成对。}同样，祂命令要聚集七套给祂，由成对组成，包括公和母——一公一母。我还能说什么？甚至不洁的鸟也不准每只带两只母鸟进入。\par

% structure: p0014 source=h3 target=section reason=relative-heading-rank; base=h3

\section{第五章 这些原始见证与基督的联系}

至此，关于原初事物的见证、我们起源的认可、以及神圣建立的预先判断，它当然是一条法律，而不仅仅是纪念，因为如果它\Quote{从起初就这样行了，}我们发现自己被基督带领到起初。正如在离婚的问题上，祂说那是\Person{梅瑟}因他们心硬而准许的，但从起初并不是这样，祂无疑将婚姻的独一性（法律）召回到\Quote{起初}。那些天主\Quote{从起初}结合在一起的，\Quote{两人成为一体}，人今日不可分开。宗徒在写给厄弗所人时也说，天主\Quote{在祂自己内预定了，在时期圆满的时分，将万有召回元首}（即起初）\Quote{一切在天上的和在地上的事物，都在基督内}。同样，希腊的两个字母，第一个和最后，主将它们归于自己，作为开始和结束的象征！它们在祂自身内汇合：为好比亚尔法一直滚动直到\OriginalTerm{奥梅加}{Omega}，而\OriginalTerm{奥梅加}{Omega}又回溯直到\OriginalTerm{亚尔法}{Alpha}，同样祂显示在祂自身内既有从开始到结束的下行过程，也有从结束回到开始的上行过程；因此每一个计划，在经由那位开始者（即天主的圣言，成了血肉的那一位）而结束于祂时，都能有一个与其开始相应的结局。所以的确，在基督内万有都被召回到\Quote{起初}，以至于信仰从割礼回到那（原始）血肉的完整，正如\Quote{它从起初就是如此；}食物的自由和唯独戒血，正如\Quote{它从起初就是如此；}婚姻的独一性，正如\Quote{它从起初就是如此；}离婚的约束，它\Emph{不是}\Quote{从起初；}最后，整个人回到乐园，在那里他\Quote{从起初}就在那里。那么，祂为何不该至少将\Person{亚当}以一夫一妻者的身份修复于彼处呢？既然祂不能以他被逐出时那样的完全完美来呈现他？因此，就起初的恢复而言，你所生活的安排以及你的希望的逻辑都要求你，凡\Quote{从起初}的，都应合乎\Quote{起初；}那个（起初）你在\Person{亚当}身上算起，在\Person{诺厄}身上重述。作出你的选择，你将哪一个视为你的\Quote{起初}。在两者中，一夫一妻制的评判权都要求你归属它。但再说一次：如果开始通向结束（如\OriginalTerm{亚尔法}{Alpha}到\OriginalTerm{奥梅加}{Omega}），正如结束回到开始（如\OriginalTerm{奥梅加}{Omega}到\OriginalTerm{亚尔法}{Alpha}），因此我们的起源被转移到基督，属血肉的转化为属神的——因为\Quote{那属神的不是在先，而是那属血肉的；然后才是属神的}——让我们照样（如之前）来看，你是否也将这同一件事归于此第二个起源：是否最后的亚当也以与第一个相同的形貌与你相遇；因为最后的亚当（即\Person{基督}）是完全未婚的，正如第一个亚当在被放逐之前也是未婚的。但是，为了你们软弱，祂献上祂自己血肉的榜样作为礼物，更完美的亚当——即\Person{基督}，在此（及其他方面）更完美，因为祂更完全纯洁——站在你面前，如果你愿意（效法祂），作为一个自愿在肉身中保持独身者。然而，如果你达不到（那种完美），祂作为一夫一妻者站在你面前，在精神上，以唯一教会为配偶，按照\Person{亚当}和\Person{厄娃}的预象，宗徒将此预象解释为基督与教会那伟大的圣事，（教导说）通过属神的，它与属血肉的一夫一妻制相类似。因此，你看，你如何在基督内更新你的起源，你不能不宣认一夫一妻制而追溯那个（起源）；除非，也就是说，你在肉身中成为祂在精神中所是的；尽管如此，祂在肉身中所是的，你也同样应当曾是。\par

% structure: p0016 source=h3 target=section reason=relative-heading-rank; base=h3

\section{第六章 \Person{亚巴郎}的案例及其对当前问题的意义}

但是，让我们继续探究我们起源的几位杰出始祖：因为有些人看来，我们一夫一妻的始祖亚当和诺厄并不令他们喜悦，或许基督也不。最后，他们诉诸亚巴郎；尽管他们被禁止承认除天主以外的任何父亲。姑且承认亚巴郎是我们的父亲；也承认保禄是我们的父亲。他说：\BibleQuote{在福音中，我生了你们。}显明你甚至是亚巴郎的儿子。因为你的起源在于他，你必须知道，这并不关涉他生命的每个时期：有一个确定的时间他是你的父亲。因为如果\Emph{信德}是源头，我们因此被算作亚巴郎的\Emph{儿子}（如宗徒教导的，对迦拉达人说的：\BibleQuote{你们知道，由此可见，那些出于信德的人，正是亚巴郎的子孙}），那么亚巴郎\BibleQuote{相信了天主，这就算作他的义德}是在\Emph{什么时候}？我认为是在他还是一夫一妻的时候，因为那时他尚未受割损。但若他后来改变了任一（相反情况）——通过与婢女同房而再婚，以及通过盟约的印记而受割损——你就不能承认他是你的父亲，除非是在他\BibleQuote{相信天主}的那个时候，如果你真是按照\Emph{信德}而非血肉成为他的儿子的话。否则，如果你跟随后来的亚巴郎作为你的父亲——即再婚的亚巴郎——那么你也当接受他的割损。如果你拒绝他的割损，那么你也会拒绝他的再婚。他的两个特质在两个方面各不相同，你无法将它们混合。他的再婚始于割损，他的一夫一妻始于未受割损。你接受再婚，也当接受割损。你保留未受割损，你也必须持守一夫一妻。此外，就像你是未受割损的亚巴郎的儿子一样，你也是那未受割损的信德之子，这是如此真实：如果你受了割损，你就立刻不再是他的儿子，因为你将不是\Quote{出于信德}，而是出于那曾在未受割损中被称义的信德的\Emph{印记}。你有宗徒：与他一同学习，连同迦拉达人。同样，如果你使自己陷入再婚，你就不是那个\Emph{信德}先于在一夫一妻中的亚巴郎的儿子。因为虽然后来他被称为\BibleQuote{万民之父}，但那是\Emph{那些}作为先于再婚的\Emph{信德}果实的民族，他们要被算作\Quote{亚巴郎的子孙}。\par

从此以后，让事情自行其道。预象是一回事，法律是另一回事；形象是一回事，律令是另一回事。形象在应验时消逝，律令则永远有待应验。形象预言，律令治理。亚巴郎那次再婚预示什么，同一位宗徒已充分教导，他是两约的解释者，正如他也规定，我们的\BibleQuote{后裔是在依撒格内被召叫的}。如果你是\BibleQuote{出于自由妇}，属于依撒格，那么他至少保持了婚姻的合一至终。\par

据此，我想，这些人便是我的起源所在。其他所有人我都不予考虑。如果我环顾他们的榜样——有些达味用甚至流血的手段为自己积聚婚姻，有些撒罗满拥有众多妻子如同财富——你被命令\Quote{追求更好的事物}；而且你还拥有若瑟，他仅结婚一次，就此我敢说比他的父亲更好；你还有梅瑟，天主的亲密目击者；你还有大司祭亚郎。第二个民族的第二个梅瑟，他引领我们的代表进入天主的恩许（之地），在他身上，那（耶稣的）圣名首次被启用，他不是再婚者。\par

% structure: p0020 source=h3 target=section reason=relative-heading-rank; base=h3

\section{第七章 从祖辈特土良转向律法先例}

在圣祖们的古代例证之后，我们同样要转向法律圣经的古代文献，以便按顺序处理我们全部的正典。既然有些人有时宣称他们与法律无关（法律基督并未废除，而是满全了），有时又随意抓住法律中的某些部分；我们也明确断言，法律在以下意义上已经过时：它的负担——依照宗徒们的判决，连祖先们也承担不了——已经完全停止；然而，那些关于正义的部分不仅永久保留，甚至被扩展了；当然，为使我们的正义能够超越经师和法利塞人的正义。如果必须要有\OriginalTerm{义德}{righteousness}，当然也必须有\OriginalTerm{贞洁}{chastity}。那么，既然法律中有这样一条诫命：若一个人没有子女就死了，他的兄弟要娶他的妻子，为给兄弟立后；而且这种情况可能重复发生在同一个人身上，正如撒杜塞人那诡诈的问题所表明的；人们因此认为在其他情况下也允许多次婚姻：那他们就应先理解这条诫命本身的原因；这样他们就会知道，那原因如今已停止，属于法律中被废除的部分。兄弟若无子女而死，必须有人继承他的婚姻：首先，因为那古老的祝福\BibleQuote{你们要生长繁衍}仍在进行中；其次，因为祖先的罪孽连子女也要承担；第三，因为阉人和不生育的人被视为可耻的。因此，唯恐那些并非因天然无能，而是因过早死亡而无子女的人，被同样视为受诅咒的；为此，一种替代的、可以说是死后所生的后裔被提供给他们。但是如今，当\BibleQuote{时日的末刻}取消了\BibleQuote{你们要生长繁衍}的命令时，因为宗徒们（另一命令）\BibleQuote{那剩下的，就是有妻子的要像没有一样}，因为\BibleQuote{时间是短促的}；而\BibleQuote{父亲}所嚼的\BibleQuote{酸葡萄}已不再\BibleQuote{使儿子的牙齿酸倒}，因为\BibleQuote{各人必因自己的罪恶而死}；而且\BibleQuote{阉人}不仅不再蒙羞，甚至配得恩宠，被邀请进入\BibleQuote{诸天的国度}：因此，继承兄弟妻子的法律被埋葬后，其反面得以确立——即\Emph{不}继承兄弟妻子的法律。这样，正如我们之前所说，当理由消失后不再有效的东西，不能为其他事情提供论证依据。所以，妻子在丈夫死后不应再婚；因为如果她再婚，她当然就是在嫁给他的兄弟：因为\BibleQuote{我们全是弟兄}。再者，女人若打算结婚，必须在\BibleQuote{主内}结婚；也就是说，不是与外邦人，而是与弟兄，因为古法律也禁止与外族通婚。此外，在肋未纪中也有警告：\BibleQuote{谁若娶了他兄弟的妻子，这是不洁——淫秽；他必无子女而死}；毫无疑问，当男人被禁止再婚时，女人也被禁止，因为她除了兄弟之外无人可嫁。那么，宗徒与法律（他并非在整体上攻击它）之间如何达成一致，等我们论到他自己的书信时再说明。暂时，就法律而言，从中引出的论证更适合我们（而非我们的对手）。简而言之，同一法律禁止司铎再婚。它还命令，司铎的女儿若成为寡妇或被休弃，且没有后代，应返回父家，由父亲的饼养育。其所以\BibleQuote{她若没有后代}，并不是说若有就可再婚——因为她若有儿子，岂不更会避免再婚？——而是说，她若有，应由儿子\Quote{养育}而非父亲；这样儿子也履行天主的诫命\BibleQuote{孝敬父亲和母亲}。此外，耶稣，圣父的至高和大司祭，用他自己的财富披戴我们——因为\BibleQuote{凡在基督内受洗的，就是穿上了基督}——使我们成为\BibleQuote{天主圣父的司祭}，按照若望所说。祂之所以召回那急着去埋葬父亲年轻人，是为了表明我们是祂所召叫的司铎；法律曾禁止司铎参加父母的葬礼：\BibleQuote{司祭不可进入任何死者的灵魂，也不可因自己的父母而沾染不洁}。\Quote{这样我们也要遵守这一禁令吗？}\Quote{当然不。}因为我们唯一的父亲——天主——\Emph{活着}，我们的母亲——教会——也活着；我们这些为天主而活的人并未死，我们也不埋葬我们的死者，因为他们也在基督内活着。无论如何，我们是基督所召叫的司铎；我们是一夫一妻制的债务人，根据天主最初的法律，那法律当时藉其司铎预言了我们。\par

% structure: p0022 source=h3 target=section reason=relative-heading-rank; base=h3

\section{第八章 从法律到福音：\Person{戴尔都良}先举出例证，再讨论教义}

现在转向那属于我们的法律——即福音——，我们遇到什么样的榜样，直到我们进入明确的教义？看，仿佛在门槛上，立即呈现在我们面前的是基督徒圣洁的两位女祭司：一夫一妻制与节欲；一位谦逊，表现在司祭\Person{匝加利亚}身上；一位彻底，表现在前驱\Person{若翰}身上；一位平息天主，一位宣讲基督；一位宣告一位完美的司祭，一位展示\Quote{比先知还大}——那一位，即不仅宣讲或亲自指出，甚至为基督施洗。因为谁更有资格在主的身躯上施行入门礼仪，比那与孕育并诞生那（身体）的同类肉身更合适？事实上，是一位童贞女，在分娩后即将一次而永远地结婚，生了基督，为使每一种圣洁的称号都能在基督的出身中得以满全，借着一位既是童贞女又是一个丈夫的妻子的母亲。再者，当祂作为婴孩被献于\Place{圣殿}时，是谁将祂接在手中？是谁首先在心灵中认出祂？是一位\Quote{正义而虔诚}的人，当然不是再婚者（这一点甚至由此考虑明显），免得基督立刻被一个女人、一位年迈的寡妇、\Quote{一个丈夫的妻子}更相称地宣讲；这女人专心地生活在\Place{圣殿}中，已经在亲自表明，何种人应当属于这属神的\Place{圣殿}——即教会。主在婴孩时期找到了这样的见证人；在成年时期，祂也没有不同的见证人。我只找到\Person{伯多禄}——通过提到他的\Quote{岳母}——是结过婚的。基于对教会的考虑，我倾向于推定他是一夫一妻者，因为教会建立在他身上，注定要从一夫一妻者中设立她圣秩的每一等级。其余的人，既然我没有发现他们结婚，我必须必然理解为他们是阉人或节欲者。诚然，在希腊人中，由于习俗的漫不经心，女人与妻子被归在一个共同的名称下——然而，有一个专指\Emph{妻子}的名称——我们因此就会如此解释\Person{保禄}，好像他表明宗徒们有妻子吗？因为如果他在辩论婚姻问题，正如他在下文所做的，宗徒本可以更好地举出某个具体的例子，那么他这样说似乎是适当的：\Quote{难道我们没有权柄带着\Emph{妻子}往来，如同其余的宗徒和\Person{刻法}一样吗？}但是当他接着加上那些表明他放弃要求供给的表述，说：\Quote{难道我们没有权柄吃喝吗？}他并没有证明宗徒们带着\Quote{妻子}往来，因为连没有妻子的人也有权吃喝；而只是带着\Quote{妇女}，她们曾像跟随主时那样服侍他们。但是，如果基督责备经师和法利塞人坐在\Person{梅瑟}的教席上，却不行他们所教导的，那么祂自己又怎会把那些只想着吩咐别人、却不同样实践肉身圣洁的人（而这种圣洁祂曾以各种方式推荐给他们去教导和实践）安置在祂自己的教席上呢？——首先以祂自己的榜样，然后以所有其他论据；同时祂告诉他们，\Quote{诸天的国}是\Quote{属于孩子的}；同时祂使那些在婚姻之后仍保持（或成为）童贞的人与这些（孩子）联合；同时祂召叫他们效法鸽子的纯朴，这只鸟不仅无害，而且谦逊，且一雄只知一雌；同时祂否认撒玛黎雅妇人的（伴侣）是丈夫，为表明多夫便是奸淫；同时，在祂自己荣耀的启示中，祂从众多圣人和先知中，偏爱\Person{梅瑟}和\Person{厄里亚}与祂同在——一位是一夫一妻者，另一位是自愿的独身者（因为\Person{厄里亚}无非就是\Person{若翰}，他带着\Quote{厄里亚的能力和精神}而来）；同时，那\Quote{贪食好酒的人}，\Quote{与税吏和罪人一同坐席、常赴午宴和晚宴的人}，仅一次参加了一场婚礼，尽管当然有许多人在（祂周围）结婚；因为祂只愿意\Emph{参加}（婚姻），正如祂只愿意它们\Emph{存在}一样。\par

% structure: p0024 source=h3 target=section reason=relative-heading-rank; base=h3

\section{第九章 从榜样，德尔图良转向直接教义教导。他以主的教导开始。}

但纵然可能有人认为这些论据是被强加且基于推测，若没有与它们平行的、主在论及离婚时所讲的教义教导，离婚原曾被允许，现今祂禁止之，首先是因为\BibleQuote{起初并不是这样}，如同婚姻的复数；其次，因为\BibleQuote{天主所结合的，人不可分开}——即恐怕他违抗主：因为唯有那结合的，祂将分开（分开，并非通过离婚的严苛，祂指责并约束离婚，而是通过死亡的债务），如果确实\BibleQuote{两只麻雀中的一只若不蒙天父的旨意，不会掉在地上}。因此，如果天主所结合的，人不可用离婚分开，那么同样合理的是，天主用死亡所分开的，人不可用婚姻再结合；再结合被分开的，将同样违背天主的旨意，正如分离结合的那样。\par

以上论及天主旨意之 \Emph{不}废，以及\Quote{起初}律法之 \Emph{重}建。但另有原因也与之一致；不，并非另一原因，而是那设立\Quote{起初}律法、并催动天主旨意禁止休妻的原因：即凡休妻者（除非因奸淫之故）便使她犯奸淫；而凡娶被休之妇者，当然也犯奸淫。被休的妇人甚至不能合法地再婚；若她未经婚姻之名而行此事，岂不归于奸淫之列？因为奸淫在婚姻之事上乃是罪。这便是天主的判定，比人的判定更为严格：普遍而言，无论通过婚姻还是随意结合，与第二个男子（交往）都被祂宣告为奸淫。因我们要看，在天主眼中何为婚姻；如此我们也将知道何为奸淫。婚姻即是：天主将\Quote{二人成为一体}联合；或者，发现（他们）已在同一肉身中联合时，祂给这结合盖上印鉴。奸淫即是：当二人——以任何方式——被\Emph{分离}后，别的——甚至可说是外来的——肉身与（其中一方）混合：这肉身不能被断言：\Quote{这是我骨中的骨，这是我肉中的肉。}因为这事从起初一次完成并宣告，如今亦然，不能适用于\Quote{别的}肉身。因此，你若说天主不愿将离异的妇人\Quote{在丈夫活着的时候}嫁给别人，仿佛祂愿意\Quote{在丈夫死后}如此行，那是毫无理由的；因为若她在他死后不受约束，活着时也一样不受约束。\Quote{正如离婚拆散婚姻，死亡亦然，她不再受那已断开结合之人的约束。}那么，她将受谁的约束？在天主眼中，她是在生前或死后结婚，无关紧要。因为她并非得罪他，而是得罪自己。\InnerQuote{人所犯的任何罪都在身体以外；但犯奸淫的人却是得罪自己的身体。}然而——正如我们先前所确立的——凡将\Quote{别的}肉身与那原始的肉身（天主当初或合二为一、或发现已经结合者）相结合的人，便是犯奸淫。祂废除那\Quote{不是从起初}的休妻，其理由正是为了巩固那\Quote{从起初}的——即\Quote{二人成为一体}的永久结合；以免第\Emph{三}种肉身的结合之必要或机会闯入（祂的领域）；除了唯一情况——即（若）那防范的（恶事）已事先发生——之外，不准许休妻。此外，休妻\Quote{不是从起初}是如此确实，以致在罗马人中，直到建城后第六百年，这种\Quote{硬心}之事才被记载为发生过。但\Emph{他们}即使不休妻也纵情于随意奸淫；而\Emph{我们}即使休妻，再婚也是不合法的。\par

% structure: p0027 source=h3 target=section reason=relative-heading-rank; base=h3

\section{第十章 圣保禄关于此主题的教导}

从这一点我看出，我们被呼吁依从宗徒的教导；为了更容易理解其含义，我们必须更热切地强调（这个论断）：妻子在丈夫死后更应约束自己，不得接受（婚姻）另一位丈夫。因为让我们思考：离婚或是因不和谐而起，或引起不和谐；而死却是天主法律的结果，而非人的过错；且是一笔所有人（甚至未婚者）都亏欠的债。因此，如果一位离婚的妇人，因不和谐、愤怒、仇恨及其原因——伤害、侮辱或任何抱怨理由——而在灵魂和肉体上都与（她的丈夫）分离，她对他——一位个人敌人，更不用说丈夫了——负有义务；那么，一位既非因自己亦非因丈夫的过错，而是因主的法律的结果，被她的配偶——不是分离，而是遗留在身后——的人，她对他（即使死了）更是负有义务，因为即使死了，她也欠他和谐？从她未曾听到（任何）离婚（话语）的人那里，她不会转身离开；她与他同在，她不曾给他写过任何离婚（文件）；她不愿失去的人，她保住了。她在内心拥有思想的自由，它向一个人以想象的享乐呈现他所没有的一切。简而言之，我询问妇人自己：\Quote{告诉我，姊妹，你让你的丈夫（安息）在平安中先你而去了吗？} 她会回答什么？（她会说）\Quote{在不和谐中？} 那样她就更负有义务，因为她与他在天主的法庭上有一桩（案子）要申诉。那位（对另一个）负有义务的人并未离开。但（她会说）\Quote{在平安中？} 那样，她必须坚持那平安，与她将不再有权离婚；并不是说，即使她能够离婚，她就可以再婚。事实上，她为他的灵魂祈祷，并为他祈求暂时的安息，以及在初次复活中的（与他）共融；并在他的安眠周年日献上（她的祭献）。因为，除非她做这些事，否则实际上就真实地离婚了，就她所能而言；而且更不义——因为她是\Emph{尽己所能}（做了离婚的事），正是由于她\Emph{没有}能力（这样做）；并且更侮辱人，因为更侮辱的是（她这样做是因为）他不\Emph{应}得到。或者，难道我们死后就停止存在，按照某个\Person{伊壁鸠鲁}（的教导），而不是按照基督的？但如果相信死者的复活，我们当然对他们负有义务，我们与他们注定要复活，互相交代。\Quote{但如果\BibleQuote{在那一世，他们不娶也不嫁，却像天使一样}，没有了婚姻关系的恢复，岂不是我们不\Emph{必}对我们的亡偶负责吗？} 不，相反我们更\Emph{要}负责，因为我们注定要进入更好的状态——注定要上升到灵性的伴侣关系，认识我们自己以及属于我们的人。否则我们怎能永远向天主歌唱感恩，如果在我们里面不再存有对这笔债的感觉和记忆？如果我们在实质上\Emph{重新}塑造，而不是在意识上？因此，我们这些将与天主同在的人将在一起；因为我们都将与同一天主同在——尽管报酬各不相同，尽管在同一位父的家里有\BibleQuote{许多住处}，为同样的雇佣的\BibleQuote{一个德纳}工作，即永生；在永生中，天主更不会分开祂所结合的，正如在今生祂禁止他们分离。\par

既然这样，一个女人如何能容纳另一位丈夫，而她即使在将来也属于她自己的（丈夫）？（此外，我们是对两性说话，尽管我们的论述只针对一方；因为相同的规定适用于双方。）她将有一人在灵性上，一人在肉体上。这将是一种奸淫：一个女人对两个男人的有意识的感情。如果其中一人已与她的肉体分开，但仍留在她心里——在那种地方，即使没有肉体接触，思想也能通过贪欲预先实行奸淫，通过意愿预先实行婚姻——他至今仍是她的丈夫，拥有使他成为丈夫的那个东西，即她的心；如果另一个人在那里找到居所，这将是犯罪。此外，如果他已从较低下的肉体交往中退出，他并未被排除。他越是纯洁，就越尊贵。\par

% structure: p0030 source=h3 target=section reason=relative-heading-rank; base=h3

\section{第十一章 关于\Person{圣保禄}教导的进一步评述}

如今，且假设你按照法律和宗徒的教导\BibleQuote{在主内}结婚——如果你尚且在乎这事——你怎有脸面向那些人请求（举行）一桩婚姻，而这对你所请求的他们而言是非法的；向一位遵循一夫一妻制的主教、向受同一庄严誓约约束的司铎和执事、向那些你亲身拒绝其修会的寡妇请求？他们明摆着会像施舍面包屑一样赐给你丈夫和妻子；因为这就是他们对\BibleQuote{凡求你的，就给他！}的解释。他们将在一位贞洁的教会——唯一基督的唯一新妇——中把你结合。而你将为你的\Emph{丈夫们}祈祷，新的和旧的。你选择吧，你要对两者中的哪一个扮演奸妇的角色。我想，两者都。但如果你有任何智慧，就为死者保持沉默。让你的沉默成为对他的离婚，另一人的嫁妆已经对此作了背书。这样，如果你忘记旧人，你就会赢得新丈夫的欢心。你应当更加努力取悦他，因为你为了他的缘故而不愿取悦天主！\OriginalTerm{属血气者}{Psychics}会认为宗徒认可了这样的（行为），或者根本没有想到，当他写道：\BibleQuote{丈夫活着的时候，妻子是被束缚的；如果丈夫死了，她就自由了，可以随意嫁人，只要是在主内。}因为正是从这段话中，他们为再婚的自由辩护，甚至任意多次（再婚），如果第二次（婚姻）可以的话：因为那已不再\Emph{是一次}的，就对\Emph{任何次数}都开放了。但宗徒究竟以何种意义写下这段话，只有先一致认定他\Emph{没有}以属血气者所利用的意义写下它，才会显明。此外，如果能先回想那些与所论段落不同的（段落），并以教义、意志和保禄自己的纪律为标准加以检验，就能达成这样的一致。因为，如果他说许第二次婚姻，而这\Emph{并不是}\BibleQuote{从起初}就有的，那他怎能断言在基督内万有都被收归于起初呢？如果他愿意我们重复婚姻关系，那他怎能主张\BibleQuote{我们的后裔}是以只结过一次婚的\Person{依撒格}为始祖而被呼召的呢？如果他使一夫一妻制成为他安排整个教会秩序的基础，那么这条规则若不先在教友——教会秩序从其等级中产生——身上生效，他又怎能如此安排呢？如果他告诫仍在婚姻中的人说\BibleQuote{时候已经缩短了}，让他们远离婚姻的享受，那他又怎能把那些借着死亡逃脱了婚姻的人重新召回婚姻呢？如果这些（段落）与当前所讨论的那一段不同，那么（如我们所说）就会一致公认：他没有以属血气者所利用的意义写下这段话；因为更易相信的是：那一段话有某种与其他段落相符的解释，而非一位宗徒似乎教导了相互矛盾的原则。我们将在主题本身中发现这种解释。是什么主题促使宗徒写下这样的（话）？一个新生的、刚刚兴起的教会的稚嫩经验，他正在喂养她，即\BibleQuote{用奶}，尚未用更强健教义的\BibleQuote{坚实的食物}；如此稚嫩，以至于信仰的婴孩状态阻止他们知道当如何处理肉身和性的需要。这种（稚嫩）的各个阶段本身从宗徒的回函中便可理解，当他说：\BibleQuote{关于你们信上所写的事，我认为男人不亲近女人倒好；但因着淫乱，男人当各有自己的妻子。}他表明有人处于婚姻状态中\BibleQuote{被信仰所擒获}，担心此后他们享受婚姻可能不合法，因为他们相信了基督的圣肉。然而他是\BibleQuote{出于让步}做出许可，\BibleQuote{不是出于命令}；即，宽容而非命令这种做法。另一方面，他\BibleQuote{宁愿}众人都像他自己一样。同样，在发出关于休妻的回函时，他表明有些人也在思考这事，主要是因为他们也认为自己信德后不应继续与外邦人的婚姻。他们进一步寻求关于\BibleQuote{童贞者}的劝告——因为\BibleQuote{主的命令}是没有的——（被告知）\BibleQuote{我认为人如果这样永久不变是好的}；（\BibleQuote{这样}），当然，如同一个人可能被信仰所见时的状态。\BibleQuote{你有妻子束缚着吗？不要寻求解脱；你没有妻子束缚着吗？不要寻求妻室。}\BibleQuote{你若娶妻，并不是犯罪}；因为对一个在相信之前就\BibleQuote{从妻子得解脱}的人，她在相信之后成为的第一个妻子，将不被算作\Emph{第二个}：因为我们的生命本身是从我们相信的时刻才开始起源的。但在这里他说他\BibleQuote{宽容他们}；否则\BibleQuote{肉身的压力}很快就会随之而来，由于时势的紧迫，那时刻躲避婚姻的累赘：是的，更应当关心如何赢得主的喜悦而非丈夫的喜悦。于是他收回了他的许可。所以，正是在同一段话中，他明确裁定\BibleQuote{各人应当永远安于他所蒙召的身份}；接着说：\BibleQuote{丈夫活着的时候，妻子是被束缚的；如果丈夫安眠了，她就自由了：她愿意嫁谁就嫁谁，只要是在主内}，他因此也表明：这样的妇人应被理解为她自己也凭着信仰\BibleQuote{被找到}\BibleQuote{从丈夫得解脱}，同样丈夫\BibleQuote{从妻子得解脱}——\BibleQuote{解脱}当然是通过死亡，而非休妻；因为对于被\Emph{休}的，他不会允许再婚，与起初的诫命相悖。所以\BibleQuote{妇人如果结婚，并不犯罪}；因为在她相信之后成为第一个丈夫的人将不被算作第二个丈夫，（同样如此娶的妻子也不被算作第二个妻子）。事实确是如此，他\Emph{因此}加上\BibleQuote{只要是在主内}；因为所争论的是关于那个曾有一个\Emph{外邦人}（丈夫）、并且在\Emph{失去他之后}才相信的妇人：唯恐她会以为甚至在相信之后还能嫁一个外邦人；尽管连\Emph{这点}也不是属血气者所关心的。让我们清楚知道，在希腊原文中，它并不以（通过两个音节的狡猾或简单改动）而成为通行的那种形式\BibleQuote{但是如果她的丈夫\Emph{将要}安眠}，好像是在说将来，因而似乎属于那位在已处于相信状态时失去丈夫的妇人。如果真如此，不加限制的许可就会允许在失去一个丈夫后任意多次地再得（新）丈夫，而没有任何甚至对外邦人也合适的婚姻节制。但即便它是这样，如同指将来时，如果某位（妇女的）丈夫\Emph{将要}死去，甚至将来同样属于那位在相信之前丈夫就死去的妇人。你任意理解吧，只要不推翻其余部分。因为既然这些（其他段落）与（上述）意义一致：\BibleQuote{你被召为奴隶，不要介意}；\BibleQuote{你是在未受割损蒙召的，不要受割损}；\BibleQuote{你是在受割损蒙召的，不要变成未受割损}；与此相符的是，\BibleQuote{你有妻子束缚着吗？不要寻求解脱；你没有妻子束缚着吗？不要寻求妻室}——显然足够表明，这些段落属于那些发现自己在一个新的、近期的\BibleQuote{蒙召}中，并就他们被信仰\BibleQuote{擒获}时所在的那些（处境状况）咨询（宗徒）的人。\par

这将是那一段落的解释，需检验它是否合乎时期与场合，以及是否与前文的事例和论据、后文的语句与含义一致，且首要地是否与宗徒本人的个別劝告和实践\OriginalTerm{实践}{practice}相符：因为没有什么比（确保）无人自相矛盾更需谨防的了。\par

% structure: p0033 source=h3 target=section reason=relative-heading-rank; base=h3

\section{第十二章　论属血气者所提出的经文解释}

此外，请听对方极为精微的论辩。他们说：\Quote{宗徒确实准许了多次婚姻，以致只有属于圣秩\OriginalTerm{圣秩}{Clerical Order}的人他才严格约束于独婚\OriginalTerm{独婚}{monogamy}之下；因为他给某些人规定的，并未给所有人规定。}那么，是否因此只对主教们\Emph{不}规定他给所有人都命令的事？——如果他给主教们规定的，并\Emph{未}给所有人命令？还是说，正\Emph{因为}给主教们，\Emph{所以}给所有人？也正\Emph{因为}给所有人，\Emph{所以}给主教们？因为主教和圣职人员从何而来？岂不是从\Emph{所有人}中来？如果\Emph{所有人}都不受独婚的约束，那么从何处选取独婚者进入圣职阶层？难道要设立一个单独之独婚者等级，用以选充圣职团体吗？（不）；但当我们自高自大、反对圣职人员时，我们就说\Quote{我们都是一体}；那时就说\Quote{我们全是司祭，因为祂使我们成为祂的天主和父的司祭}。当我们被要求与司祭纪律完全平等时，我们便放下（司祭的）冠带，却仍处于同一层次！宗徒写信时所论及的是关于教会圣秩——哪些人子应当被祝圣。因此，普遍纪律的全部形式理应摆在首要位置，如同一个谕令，在某种意义上普遍且应细心遵守，好使平信徒更清楚地知道，他们自己必须遵守那对他们的监督者不可或缺的秩序；甚至荣誉职位本身也不可因任何纵容而自夸，仿佛基于地位的特权。圣神预见到有些人会说：\Quote{主教们凡事都可行}；正如你们那位乌提纳的主教，甚至连斯坎提尼法也不畏惧。唉，在你们的教会中，有多少再婚者主持圣事，显然是在侮辱宗徒；至少在他们在座时读这些经文，他们并不脸红！\par

来吧，你们这些认为独婚是为主教们制定的例外法律的人，也请放弃其余那些与独婚一同归于主教的纪律名目。拒绝成为\Quote{无可指摘、节制、品行端正、有序、好客、易受教导}；相反，（成为）\Quote{好酒、好斗、好争吵、贪财、不治理自己的家、不关心子女的管教}，也不\Quote{甚至从外人那里寻求美名}。因为如果主教们有自己教导独婚的法律，那么其他那些与独婚相称的品行也就唯独为主教们所记载。但对于平信徒，既然独婚对他们不合适，其他品行也与之无关。（这样），属血气者，你（若愿意）就逃脱了整个纪律的束缚！请一致地规定：\Quote{给某些人命令的，并未给所有人命令}；否则，如果其他品行确实是共通的，而独婚却只加给主教们，请问：难道只有\Emph{他们}，承受全部纪律的人，才能被称为\Emph{基督徒}吗？\par

% structure: p0036 source=h3 target=section reason=relative-heading-rank; base=h3

\section{第十三章　答复来自圣保禄的进一步反对意见}

\Quote{但再写信给弟茂德时，他\BibleQuote{愿意年轻的妇女嫁人，生育儿女，治理家务。}}他（在此）是对他上面所称的那些人说话——即\Quote{年轻的寡妇}，她们在守寡后被捕获\OriginalTerm{被捕获}{apprehended}，并经过一段时间的追求，且在情感中已有了基督之后，\Quote{愿意嫁人，但受审判，因为他们废弃了起初的信德}——即她们在守寡中\Quote{被找到}，且宣认之后未能坚持的信德。因此，他\Quote{愿意}她们\Quote{嫁人}，免得她们以后废弃所宣认的守寡之初信；并非准许她们一再嫁人，只要她们拒绝在充满诱惑的——不，毋宁是沉溺于放纵的——守寡中坚持。\par

\Quote{我们读到他在\BibleBook{}{罗马书}中写道：\BibleQuote{但丈夫活着的女人，是受丈夫的法律约束；丈夫若死了，她就从丈夫的法律中解脱了。}毫无疑问，丈夫活着，她若嫁给第二个丈夫，就会被视为犯奸淫。然而，丈夫若死了，她就从（他的）法律中得了解放，因此她若作（妻子）给另一个丈夫，就不是犯奸淫。}但也请阅读后续，以免这个取悦你们的含义从你们手中溜走。\Quote{因此，}他说，\BibleQuote{我的弟兄们，你们借着基督的身体也成了死，脱离了法律，得以归于另一位，就是从死者中复活的那一位，使我们为天主结果实。因为我们还在肉身时，那借着法律而激起的罪恶情欲，在我们的肢体中运作，结出死亡的果实；但现在我们已从法律下解脱了，因为我们已经死在了那曾约束我们的法律之下，好让我们以心神的新颖而非文字的陈旧来事奉天主。}因此，如果他吩咐我们\BibleQuote{借着基督的身体成为死，脱离法律}，（即教会，它存在于心神的新颖中），而不是\BibleQuote{借着文字的陈旧}，（即法律）——他使你们脱离法律，即法律不阻止妻子在丈夫死后成为另一个丈夫的妻子——他使你们受相反条件的约束，即你们失去丈夫后不得再婚；并且，正如你们若在（你们的）丈夫死后成为第二个丈夫的妻子，本不会被视为犯奸淫——如果你们仍必须受法律约束行事——那么，由于（你们）状况的差异，他确实预见你们在丈夫死后若另嫁他人，便是犯奸淫：因为你们如今已对法律成了死，既然你们已从那法律中退出，而那法律曾准许你们这样做，现在对你们而言便不再合法了。\par

% structure: p0039 source=h3 target=section reason=relative-heading-rank; base=h3

\section{第十四章　即使\Person{圣保禄}在属血气的人所理解的意义上给予许可，那也不过如同\Person{梅瑟}允许休妻一样——是对人心硬的一种迁就}

如今，如果宗徒甚至绝对允许在（皈依）信仰之后失去配偶的人再婚，那么他这样做，就像他违背自己规则的（严格）字句而做的其他事一样，是因应时机：他为了\BibleQuote{假冒的假弟兄}而给\Person{弟茂德}行了割礼；又因犹太人的监视而领着几个\BibleQuote{剃头的人}进了圣殿——他却责斥那些愿意遵守法律的迦拉达人。但时机要求他\BibleQuote{成为一切人，为的是救一切人}；\BibleQuote{为他受生产之苦，直到基督在他们身上成形}；并\BibleQuote{像保姆一样}抚育信仰的幼小者，\BibleQuote{以宽恕的方式，而非命令的方式}教导他们——因为\Quote{宽恕}是一回事，\Quote{命令}是另一回事——因着\BibleQuote{肉身的软弱}而暂时准许再婚，正如\Person{梅瑟}因\BibleQuote{心硬}而准许休妻一样。\par

因此，我们在此也要阐明他（这段话）的补充含义。因为如果基督废除了\Person{梅瑟}所吩咐的，因为\BibleQuote{起初并不是这样}；并且——既然如此——基督不会因此被认为来自另一个权能；那么，为什么\OriginalTerm{护慰者}{Paraclete}也不能废除\Person{保禄}所给予的宽恕呢？——因为再婚\BibleQuote{起初也不是这样}——而不必因此被怀疑是外来的神，只要这补充符合天主和基督的尊荣？如果当时机（宽恕的时期）已满时，遏制\BibleQuote{心硬}符合天主和基督的尊荣，那么当\BibleQuote{时期}已经更加\BibleQuote{缩短}时，摆脱\BibleQuote{肉身的软弱}岂不是更符合天主和基督的尊荣吗？如果婚姻不可分离是合理的，那么理所当然，不再重复也是光荣的。总之，在世人的估价中，两者都被视为良好纪律的标志：一个以和谐之名，一个以节制之名。\BibleQuote{心硬}统治直到基督的时代；让\BibleQuote{肉身的软弱}（满足于）统治到护慰者的时代为止。新法律废除了离婚——它（必须）废除某些事；新预言（废除）再婚，后者也不亚于前一个婚姻的离婚。但\BibleQuote{心硬}比\BibleQuote{肉身的软弱}更容易服从基督。后者比前者更支持\Person{保禄}；的确，如果它抓住他的宽恕却拒绝他的规诫——避开他更审慎的意见和他常有的\Quote{意愿}，不让我们向宗徒献上他所更愿意的（服从）——那就是在利用他。\par

这种最无耻的\Quote{\OriginalTerm{软弱}{infirmity}}还要持续多久，与\Quote{\OriginalTerm{更美的事}{better things}}进行歼灭战？它放纵的时间（间隔）直到\OriginalTerm{圣神（护慰者）}{Paraclete}开始祂的行动为止，主将那些在祂的日子\Quote{不能承受的}事延迟到祂的来临；如今任何人不能再无法承受，因为那授予承受能力的那一位并不缺乏。我们要到何时才引用\Quote{\OriginalTerm{血肉}{flesh}}，因为主说：\BibleQuote{血肉是软弱的}\BibleRef{玛 26:41}？但祂也已预先说过：\BibleQuote{圣神是切愿的}\BibleRef{玛 26:41}，为叫圣神能战胜血肉——叫软弱的服从更强有力的。因为祂又说：\BibleQuote{谁能领受，就领受吧}\BibleRef{玛 19:12}；也就是说，那\Emph{不能}领受的就走开。那富翁\Emph{确实}走开了，他没有\Quote{领受}分施财物给穷人的诫命，被主抛弃在他自己的意见中。因此，也不会将\Quote{苛刻}归咎于基督，而是每个人自由意志恶行的原因。\BibleQuote{看，}\BibleRef{申 30:15}祂说，\BibleQuote{我把善和恶摆在你面前}\BibleRef{申 30:15}。选择那善的：如果你不能，因为你不愿意——因为祂已表明你如果愿意就能，因为祂将两者摆在你自由意志面前——你应该离开那你不愿意顺从的那一位。\par

% structure: p0043 source=h3 target=section reason=relative-heading-rank; base=h3

\section{第十五章 论以严厉指责新预言之门徒的不公；这指责更应归咎于属魂派}

因此，我们这里有什么严厉之处，如果我们与那些不承行天主旨意的人断绝（共融）？有什么异端，如果我们判断再婚为非法，近似奸淫？因为什么是奸淫，无非是非法的婚姻？宗徒曾烙印那些惯于完全禁止婚姻的人，同时也禁止食用天主所造的食物。然而，我们并不因为摒弃婚姻的重复而废除了婚姻，就如我们并不因为更频繁地斋戒而谴责食物一样。废除是一回事，规范是另一回事；订立不结婚的法律是一回事，限定结婚的次数是另一回事。坦率地说，如果那些以严厉指责我们、或认为在这件事上存在异端的人，竟如此纵容\Quote{\OriginalTerm{血肉的软弱}{infirmity of the flesh}}，以致认为必须藉频繁的婚姻来支持它；那么，为什么在另一种情况下，他们既不给予支持，也不以宽宥纵容它——即当折磨使之否认信仰的时候？因为，当然，那在战场上跌倒的软弱比那在卧室中跌倒的（软弱）更可原谅；那在拷问台上屈服的比那在新床上（屈服）的更可原谅；那向残忍屈服的比那向情欲（屈服）的更可原谅；那呻吟着被征服的比那在冲动中（被征服）的更可原谅。但前者他们逐出教会，因为它没有\BibleQuote{坚持到底}\BibleRef{玛 24:13}；后者他们却扶持，好像它确实\BibleQuote{坚持到底}\BibleRef{玛 24:13}了。提出（问题）为什么各人没有\BibleQuote{坚持到底}\BibleRef{玛 24:13}；你就会发现，那未能忍受野蛮的软弱的原因，比那（未能忍受）节操的（软弱）更高尚。然而，即使是\OriginalTerm{血催逼的}{bloodwrung}——更不用说不知羞耻的——\OriginalTerm{背离}{defection}，\Quote{血肉的软弱}也不能为之开脱！\par

% structure: p0045 source=h3 target=section reason=relative-heading-rank; base=h3

\section{第十六章 论为再婚辩护所提出的理由之薄弱}

但我摇头哂笑，当有人以\Quote{肉身的软弱}为借口来反对我们时——其实那软弱更当称作力量的巅峰。婚姻的重复是力量的事：从节制的安逸中复活到肉身的工作，需要强有力的韧带。这种\Quote{软弱}与第三次、第四次、甚至第七次婚姻相当；它每次增加力量如同增加软弱；它不再有宗徒的权威支持，而是某个赫尔摩革涅的支持——他娶的女人比他画的还多。因为在他是质料充沛；由此他猜测连灵魂也是物质的；所以他比其他人大大\Emph{没有}从天主来的圣神，甚至不再是魂人，因为他的魂质也不是来自天主的吹气！假如有人借口\Quote{匮乏}，公开使他的肉身卖淫，为生计而婚配，却忘记了不当为饮食忧虑？他有天主眷顾，就是喂养乌鸦、装扮鲜花的天主。假如他借口家中孤独？仿佛一个永远准备逃遁的男人，一个女人就能给他作伴？当然，他身边有寡妇，他有权利娶她。不只允许娶一个，而是可以娶多个。假如有人思念后裔，有如罗特妻子的眼睛一般，因此以第一次婚姻没有孩子为理由重复婚姻？基督徒竟会寻求继承人，而他从全世界被剥夺了继承权！他有\Quote{弟兄}；他有教会作母亲。除非人们相信，在基督的法庭上，也像在罗马一样，按照尤利安法行事，并以为未婚无子者不能按天主的遗嘱领受全额。让他们这样想的人，一直结婚到终点吧；在这肉身的混乱中，他们如索多玛、哈摩辣和洪水之日，被命定的世界末日所吞噬。让他们加上第三句说话：\Quote{我们吃喝\Emph{和结婚}吧，因为明天我们就要死了；}不思想那\Quote{祸哉}斥责\Quote{给怀孕和哺乳的}，在普世的\Quote{震动}中，会比在犹太一隅的蹂躏中更加严厉痛苦。让他们因重复婚姻积蓄最合时宜的末世果实——起伏的乳房，翻腾的子宫，啼哭的婴孩。让他们为假基督预备孩子，使他能比法郎更残暴地向他们发泄。他将带来杀婴的产婆。\par

% structure: p0047 source=h3 target=section reason=relative-heading-rank; base=h3

\section{第十七章 外邦人的榜样谴责这种\Quote{肉身的软弱}}

他们将显然有貌似有理的特权在基督面前辩护——那永远的\Quote{肉身的软弱！}但审判这软弱的，不再是我们的\OriginalTerm{一夫一妻的父亲}{monogamist father}依撒格，或基督的著名自愿独身者若望，或默辣黎的女儿友弟德，以及诸如此类的众多圣人榜样。外邦人常是我们未来的判官。将有一位迦太基女王站起来，审判基督徒，她身为难民，在他乡寄居，当时正创立如此强大的国家，她本可不必请求就获得王室婚姻，却唯恐经历第二次婚姻，宁可\Quote{烧}也不\Quote{嫁}。与她同坐的还有一位罗马主妇，她虽然是通过夜间的暴力认识了另一个男人，却用血洗净了肉身的污点，以亲自报复仇来维护一夫一妻制。也有为丈夫而死、宁愿死后不嫁的人。至少对偶像而言，一夫一妻制和守寡都作侍从。只有结过一次婚的妇女才在妇女福尔图娜和玛图塔母亲上挂花环。大祭司和弗拉门的妻子只结一次婚。刻瑞斯的祭司，甚至在她丈夫生前并同意下，以友好分居成为寡妇。也有以绝对节欲来审判我们的：维斯塔贞女、亚该亚的朱诺、西徐亚的狄安娜、皮提亚的阿波罗的贞女。以节欲而言，著名的埃及公牛祭司也将审判基督徒的\Quote{软弱。}羞耻吧，你这\Quote{穿上}基督的肉身！你一次结婚就足够了，你从\Quote{起初}被造，又被\Quote{终结}召回！至少回到先前的亚当吧，如果不能回到最后的亚当！他一次尝了树上的果子；一次感到了私欲偏情；一次掩盖了羞耻；一次在天主面前脸红；一次隐藏了有罪的脸色；一次从圣洁的乐园被放逐；此后一次结婚。如果你\Quote{在他内}，你就有规范；如果你已\Quote{进入基督}，你就要更好。请给我们展示第三个亚当，一个再婚者；那样你就能成为在两者之间的你所不能成为的。\par

\SourceRecord{Translated by S. Thelwall.；Ante-Nicene Fathers；Vol. 4.；Edited by Alexander Roberts, James Donaldson, and A. Cleveland Coxe.；Buffalo, NY: Christian Literature Publishing Co.,；1885.；<http://www.newadvent.org/fathers/0406.htm>.}
